\documentclass[10pt,a4paper]{amsart}
\usepackage[margin=19mm]{geometry}
\usepackage{amsmath,amssymb,mathtools}
\usepackage[scr=rsfs,cal=euler]{mathalfa}
\usepackage{xfrac}
\usepackage[unicode,hidelinks,bookmarks=false]{hyperref}
\newcommand{\ie}{i.e.\ }
\newcommand{\cf}{cf.\ }
\newcommand{\resp}{respectively\ }
\newcommand{\Subsection}{\subsection}
\providecommand{\nobreakdash}{\nobreak\hbox{-}}
\setlength{\emergencystretch}{2em}
\multlinegap0pt
\usepackage{bm}
\usepackage{bbm}
\usepackage{dsfont}
\usepackage{xspace}
\usepackage{cancel}
\usepackage{mathtools}
\usepackage{amsthm}
\makeatletter
\@ifundefined{theorem}{\newtheorem{theorem}{Theorem}}{}
\@ifundefined{lemma}{\newtheorem{lemma}[theorem]{Lemma}}{}
\makeatother
\newcommand{\al}{\alpha}
\newcommand{\apr}{{\rm AP}(\r)}
\newcommand{\ep}{\varepsilon}
\newcommand{\ft}{\widehat}
\newcommand{\la}{\lambda}
\newcommand{\n}{\mathbb{N}}
\newcommand{\p}{\varphi}
\newcommand{\wt}{\widetilde}
\newcommand{\z}{\mathbb{Z}}
\renewcommand{\d}{\mathrm{d}}
\renewcommand{\r}{\mathbb{R}}
\title[A classification of Fourier summation formulas and crystalli]{A classification of Fourier summation formulas and crystalline measures}
\author{Felipe Gon\c{c}alves}
\date{Source: arXiv:2312.11185v3 (CC BY 4.0)}
\begin{document}
\maketitle

\noindent\emph{Source and licence.} A classification of Fourier summation formulas and crystalline measures. Version arXiv:2312.11185v3 (CC BY 4.0), distributed under the Creative Commons Attribution 4.0 International licence, as recorded in the arXiv licence metadata for this exact version and shown on the abstract page https://arxiv.org/abs/2312.11185v3. Full rights notice: licenses/arxiv-2312.11185-CC-BY-4.0.txt.\par\smallskip

\noindent\textit{Editorial scope.} Counted unit: the Lemma whose source label is \texttt{lem:polyaprox} in the article (source0.tex lines 811--813, source label \texttt{lem:polyaprox}), delivered together with the article's own complete original proof of it (source0.tex lines 814--870, one \texttt{proof} environment of 57 lines). No mathematics is written, repaired or re-derived: statement and proof are the article's own text line for line. Required context retained uncounted: none. Every definition, display and label of those lines is retained unchanged. Omitted from this package and not redistributed: 1--810, 871--1689. Nothing inside a retained excerpt is omitted or altered: every retained body is delivered exactly as the article's lines are written, so no omission receipt of any kind accompanies this package. Citations: the retained excerpts carry no citation command at all, so no bibliography entry of the article is transcribed and no reference is redistributed. Reference adaptation: none was needed and none was made; every reference inside the delivered unit resolves inside it, so no unresolved reference or citation is produced. Printed slips kept literal: no printed word, symbol, hypothesis or number of the counted statement or of its proof was altered. Layout and class: the article's own class is \texttt{amsart} with packages including inputenc, amsmath, amssymb, amsthm, epsfig, hyperref, latexsym, url, euscript, color, harpoon, mathabx, tikz, fullpage; the wrapper is the lane's pinned \texttt{amsart} editorial wrapper at 10pt on A4, which restates the theorem-like environments the retained text needs and adds the article's own macro definitions that the retained text needs (\texttt{\textbackslash al}, \texttt{\textbackslash apr}, \texttt{\textbackslash d}, \texttt{\textbackslash ep}, \texttt{\textbackslash ft}, \texttt{\textbackslash la}, \texttt{\textbackslash n}, \texttt{\textbackslash p}, \texttt{\textbackslash r}, \texttt{\textbackslash wt}, \texttt{\textbackslash z}), with extra wrapper packages (bm, bbm, dsfont, xspace, cancel, mathtools). The delivered numbering is the unit's own. Assets: no retained span contains a figure environment, a graphics-inclusion command, a PDF-page inclusion, a table or a TikZ picture; no image file is copied into this package, the package declares no asset dependency and no image file is redistributed with it. Licence: version arXiv:2312.11185v3 of this article is distributed under the Creative Commons Attribution 4.0 International licence, as recorded in the arXiv licence metadata for this exact version and shown on the abstract page https://arxiv.org/abs/2312.11185v3. A full rights notice is delivered at licenses/arxiv-2312.11185-CC-BY-4.0.txt. Characters: every retained excerpt is pure ASCII, so the delivered engine, XeLaTeX with its default Latin Modern text font, reports no missing character.
\begin{lemma}\label{lem:polyaprox}
Given any $f\in \apr$ and $\ep>0$, there is a trigonometric polynomial $p$ such that $\|f-p\|_\infty<\ep$. In particular, $\apr$ is the closure in $C(\r)$ of the algebra of trigonometric polynomials. 
\end{lemma}

\begin{proof}
We can assume $\|f\|_\infty\leq 1$. For a given $\ep>0$ let $\tau_\ep$ be the set of $\ep$-translations for $f$.  First we will need some auxiliary functions. We claim that for every sufficiently large $M\geq 1$ there exists $\p_M\in C^\infty_c(\tau_\ep \cap (-M,M))$ such that $0\leq \p_M\leq 1$, $\int_\r \p_M =1$, $\int_\r |\p_M'| =O_{\ep}(1)$ and 
$$
\sum_{n\in \z} |\ft \p(n/(4M))|^{1/2} = O_{\ep}(1).
$$
Assuming this claim is true, we now finish the proof. Let
$$
f_M(x) = \int_\r f(x+t)\p_M(t)\d t
$$
and note that $\|f-f_M\|_\infty\leq \ep$ for any $M$. Also note that if $|x|\leq M$ then
$$
f_M(x)=  \int_{-2M}^{2M} f(t)\p_M(t-x)\d t. 
$$
The function $\p_M(x)$ can be identified with a $4M$-periodic  $C^\infty$-function and so we have
$$
\p_m(x)=\sum_{n\in\z} \theta_n e^{2\pi i n x/(4M)}
$$
where $\sum_{n\in \z} |\theta_n|^2 = (4M)^{-1}\int_{-2M}^{2M} |\p_M(x)|^2\d x$ and 
$$
\theta_n = (4M)^{-1}\int_{-2M}^{2M} \p_M(x)e^{-2\pi i n x/(4M)}\d x = \ft \p_M(n/(4M))/(4M).
$$
We obtain
$$
f_M(x) = \sum_{n\in\z} \wt \theta_n e^{-2\pi i n x/(4M)}  \quad \text{ for } |x|\leq M,
$$
where $\wt \theta_n = \theta_n \int_{-2M}^{2M}f(t) e^{2\pi i n t/(4M)}\d t$. Note this representation converges absolutely since $ |\wt \theta_n| \leq  4M|\theta_n| = |\ft \p_M(n/(4M))| \leq 1$, $\sum_{n\in\z} |\wt \theta_n|^{1/2} = O_{\ep}(1)$, and so, $\sum_{n\in\z} |\wt \theta_n|=O_{\ep}(1)$. Now enumerate $(\wt \theta_n)_{n\in \z}$ in decreasing order of magnitude and call this new sequence $(\al_{M,n})_{n\geq 1}$.  We obtain
$$
f_M(x) = \sum_{n\geq 1} \al_{M,n} e^{2\pi i \la_{M,n} x}  \quad \text{ for } |x|\leq M,
$$
for some $\la_{M,n}\in \frac{1}{4M}\z$. Noticing that 
$$
|\al_{M,n}|^{1/2}n \leq \sum_{j=1}^n |\al_{M,j}|^{1/2} = O_{\ep}(1),
$$
we obtain $|\al_{M,n}|\leq n^{-2}$. A standard Cantor's diagonal procedure guarantees the existence of $(\al_n)_{n\geq 1}$ such that $\sum_{n\in\z} |\al_n| = O_{\ep}(1)$ and a subsequence $M_k\to\infty$ such that $\lim_k \al_{M_k,n}=\al_n$ for all $n\geq 1$.  Let now $I\subset \n$ be the $n$'s such that $\sup_{k} |\la_{M_k,n}| < \infty$.  By further taking a subsequence of the $M_k$'s we can assume that $\lim_k \la_{M_k,n} \to \la_n$ for $n\in I$.  Note now that
$$
|\al_{M_k,n}| \leq |\ft \p_M(-\la_{M_k,n})|  \leq  \frac{1}{2\pi |\la_{M_k,n}|} \int_{\r} |\p_M'(x)|\d x  = O_\delta(|\la_{M_k,n}|^{-1}).
$$
In particular, $\lim_k \al_{M_k,n} = \al_n=0$ if $n\notin I$. Finally, the uniform bound $|\al_{M,n}|\leq n^{-2}$ forces $f_{M_k}$ to converge uniformly in compact sets to 
$$
g(x) =  \sum_{n\geq 1} \al_{n} e^{2\pi i \la_{n} x},
$$
which in particular implies that
$$
\|f-g\|_\infty \leq \ep.
$$
However, we can now simply truncated $g$ to find a trigonometric polynomial $p$ such that $\|f-p\|_\infty \leq 2\ep$.

It remains to construct the auxiliary functions $\p_M$. Since $f$ is also uniformly continuous, it is easy to show that $\{t_n\}_{n\in \z} + (-\delta,\delta) \subset \tau_{\ep/4}$ for some small $\delta=\delta_\ep>0$ and some sequence $\{t_n\}_{n\in \z}$ satisfying $1/\delta \leq t_{n+1}-t_n \leq 3/\delta$ for all $n$. Take $h \in C^\infty_c(-1,1)$ even with $0\leq h\leq 1$ and $\int_\r h=1$. Let $h_\delta(x)=h(x/\delta)/\delta$, $\psi(x)=(2N+1)^{-1}\sum_{n=-N}^N h_\delta(x-t_n)$ and $\p_M(x)=\psi * \psi *\psi * \psi$, where $N\geq 1$ is selected to be the largest such that $\{t_n\}_{|n|\leq N} + (-\delta,\delta) \subset \tau_{\ep/4}\cap (-M/4,M/4)$. Note we must have $N\geq \frac{\delta M}{10}$ for sufficiently large $M$. Since $\tau_{\ep/4} + \tau_{\ep/4}+\tau_{\ep/4}+\tau_{\ep/4}\subset \tau_\ep$ we conclude that $\p_M \in C^\infty_c(\tau_\ep \cap (-M,M))$. Also note that $\int_\r \p_M = (\int_\r \psi)^4 = 1$ and, since $\p_M'=\psi'*\psi*\psi *\psi$, that 
$$
\int_\r |\p'_M| \leq \int _\r|\psi'|  = O(1/\delta).
$$
Finally note that 
\begin{align*}
\sum_{n\in\z} |\ft \p_M(n/(4M))|^{1/2} = \sum_{n\in\z} |\ft \psi(n/(4M))|^{2} & = 4M \int_{-2M}^{2M}|\psi(x)|^2\d x \\ & = \frac{4M}{\delta (2N+1)} \int_{-1}^1 |h(x)|^2 \d x  \leq \frac{4M}{\delta (2N+1)}  \ll \delta^{-2},
\end{align*}
for $M$ large enough. This finishes the proof.
\end{proof}

\end{document}
